\question
Вы работаете бэкенд-разработчиком в крупной компании, и недавно к вам подошел аналитик с необычным требованием: вам необходимо научиться сохранять используемые в приложении графы в базу данных (речь идет о реляционных БД, наподобие PostgreSQL). Подумав, вы поняли, что любую таблицу в БД с двумя столбцами можно описать как бинарное отношение, а бинарное отношение, в свою очередь, можно представить в виде графа. Понимая, что ситуации могут быть разными, вы решили потренироваться в построении бинарных отношений на разных типах графов. Постройте графы и эквивалентные бинарные отношения (графы) в соответствии с типом связности графа:

\begin{itemize}
\item ориентированный граф, обладающий 6 компонентами сильной связности, в которых более 3 вершин
\item неориентированный граф, состоящий из 6 компонент связности, причем каждая из них имеет не менее 3 компонент реберной или вершинной двусвязности (только для связности строим отношение)
\item неориентированный граф, состоящий одновременно из 7 компонент вершинной и 9 компонент реберной двусвязности (для обеих двусвязностей построить бинарное отношение)
\end{itemize}
ДОПОЛНИТЕЛЬНО: покажите на множестве вершин получившиеся классы эквивалентности по каждой связности!

---------------

Автор -- Тимур Гонтарь, M3308